\section{Théorème de confluence des réponses dimensionnelles}
\label{rm:ch:confluencia-general}

\subsection{Antécédent commun et produit fibré}

Les résultats du volume proviennent de lecteurs hétérogènes appliqués à un même état généalogique enrichi. Soit

\begin{equation}
 X=\bigl(X_{\APP},X_{\TRIT},X_{\TPK},
          \mathcal C,\mathcal B,\mathcal R,\mathcal M,\partial X\bigr)
 \label{rm:eq:confluence-state}
\end{equation}

l’état qui conserve les cylindres, la retenue, le parcours, la mémoire et la frontière. Sur \(X\) sont construits les lecteurs typés

\begin{equation}
 P_i:X\longrightarrow Y_i,
 \qquad
 i\in I:=\{\alpha,{\rm act},{\rm ang},108,\square,{\rm obs},{\rm Per}\}.
 \label{rm:eq:confluence-objects}
\end{equation}

Leurs codomaines représentent respectivement le générateur dodécaphasique de couplage, le couple d’action, l’opérateur angulaire bicouche, l’horloge \(108\), l’écran cubique, l’observateur interne et le mode de Perron. Pour chaque lecteur, on considère son graphe, muni de la projection canonique

\begin{equation}
 \Gamma_i:=\{(x,y_i)\in X\times Y_i:P_i(x)=y_i\},
 \qquad
 \pi_i:\Gamma_i\longrightarrow X,
 \quad (x,y_i)\longmapsto x.
 \label{rm:eq:confluence-reader-graphs}
\end{equation}

L’objet maître est le produit fibré bien typé de ces graphes :

\begin{equation}
 \boxed{
 \mathfrak M_X=
 \mathop{\times}_{X,\,i\in I}\Gamma_i
 =\left\{\bigl(x,(y_i)_{i\in I}\bigr):
 P_i(x)=y_i\ \text{pour tout }i\in I\right\}.}
 \label{rm:eq:master-fiber-product}
\end{equation}

Les projections \(\pi_i\) fixent matériellement le même antécédent. Un élément de \(\mathfrak M_X\) est une famille de sorties compatibles accompagnées de l’état qui les engendre.

\begin{definition}[Publication dimensionnelle]
Une publication dimensionnelle est une composition
\begin{equation}
 \mathfrak M_X\xrightarrow{\;P\;}\mathcal L
 \xrightarrow{\;\operatorname{ev}_{\mathcal U}\;}\RR,
 \label{rm:eq:dimensional-publication}
\end{equation}
où \(\mathcal L\) est une droite dimensionnelle typée et \(\operatorname{ev}_{\mathcal U}\) est l’évaluation finale dans une carte d’unités. L’application \(P\) construit d’abord l’invariant et sa généalogie.
\end{definition}

\subsection{Indépendance horizontale et compatibilité verticale}

Deux lecteurs issus d’un même antécédent peuvent satisfaire des relations croisées sans se factoriser l’un par l’autre. Cette distinction régit les doubles feuillets du vide, les réalisations positive et orientée, la double projection archimédienne et exceptionnelle, ainsi que les lectures géométrique, informationnelle et de saveur de la transition \(6\to8\).

\begin{theorem}[Confluence sans factorisation horizontale]
Soient \(P_i:\mathfrak M_X\to Y_i\), \(i=1,\ldots,r\), des publications construites sur le même état. Supposons que chaque \(P_i\) conserve ses données typées et que les relations \(R_j(P_1,\ldots,P_r)=0\) soient démontrées après la construction des lecteurs. Alors le morphisme conjoint
\begin{equation}
 P=(P_1,\ldots,P_r):\mathfrak M_X\longrightarrow
 Y_1\times\cdots\times Y_r
 \label{rm:eq:joint-publication}
\end{equation}
est verticalement compatible avec \(X\). Une factorisation \(P_i=F_{ij}\circ P_j\) exige en outre que \(P_j\) sépare toutes les fibres que \(P_i\) distingue. La seule existence de \(R_j=0\) n’implique pas cette séparation.
\end{theorem}

\begin{proof}
La compatibilité verticale résulte de la propriété universelle du produit fibré des graphes \(\Gamma_i\) : toutes les composantes ont la même image dans \(X\). Si \(P_i=F_{ij}\circ P_j\), alors \(P_j(x)=P_j(x')\) impose \(P_i(x)=P_i(x')\). Par conséquent, toute paire \(x,x'\) ayant la même lecture \(j\) et des lectures \(i\) différentes réfute la factorisation. Une relation scalaire \(R_j=0\) ne fait que restreindre l’image conjointe ; elle ne garantit pas l’inclusion des noyaux nécessaire à la factorisation.
\end{proof}

Le théorème explique pourquoi le produit des feuillets peut fixer le cône causal tandis que leur quotient conserve l’orientation, et pourquoi l’aire totale peut coexister avec un Hamiltonien modulaire qui distingue la provenance \(90/120\).

\subsection{1re confluence : couplage, vide et secteur électrofaible}

Soit

\begin{equation}
 D=D_A=\det T,
 \qquad
 Q(D)=-\frac9{25}\log D=4\pi\alpha_{\HMT}.
 \label{rm:eq:confluence-q-d}
\end{equation}

La mémoire orientée s’exprime par

\begin{equation}
 \rho=e^{2C^*_{\rm rad}},
 \qquad
 q_-=\sqrt{D\rho},
 \qquad
 q_+=\sqrt{D/\rho}.
 \label{rm:eq:confluence-qpm}
\end{equation}

Le calcul fonctionnel

\begin{equation}
 F(z)=\frac{1-z^{90}}{1-z^{120}},
 \qquad
 r_\pm=F(q_\pm)
 \label{rm:eq:confluence-F}
\end{equation}

produit

\begin{equation}
 \widehat\mu=r_-^2,
 \quad
 \widehat\varepsilon=r_+^2,
 \quad
 \widehat Z=\frac{r_-}{r_+},
 \quad
 \widehat c=(r_-r_+)^{-1}.
 \label{rm:eq:confluence-vacuum}
\end{equation}

À l’interface électrofaible au niveau des arbres,

\begin{equation}
 g^2=\frac{Q(D)}{1-D},
 \qquad
 g'^2=\frac{Q(D)}{D},
 \label{rm:eq:confluence-gauge-couplings}
\end{equation}

d'où

\begin{equation}
 \boxed{
 (1-D)g^2=Dg'^2=Q(D),
 \qquad
 \frac1{g^2}+\frac1{g'^2}=\frac1{Q(D)}.}
 \label{rm:eq:confluence-ew-master}
\end{equation}

La première égalité exprime le même couplage dans deux orientations de jauge ; la seconde le retrouve comme somme de compliances. Le déterminant \(D\) régit le mélange pair et \(\rho\) conserve l’orientation impaire.

La coordonnée de Higgs intègre cette dernière :

\begin{equation}
 \lambda_H=\frac{Q(D)}{8(1-D)}D^{-3}\rho^{5/3}.
 \label{rm:eq:confluence-higgs}
\end{equation}

En définissant

\begin{equation}
 \Xi=\frac{8(1-D)D^3\lambda_H}{Q(D)},
 \label{rm:eq:confluence-xi}
\end{equation}

est reconstruit

\begin{equation}
 \rho=\Xi^{3/5},
 \qquad
 \widehat Z=
 \frac{F(\sqrt D\,\Xi^{3/10})}
      {F(\sqrt D\,\Xi^{-3/10})}.
 \label{rm:eq:confluence-vacuum-higgs}
\end{equation}

L’équation \cref{rm:eq:confluence-vacuum-higgs} constitue un contrôle croisé entre orientation scalaire et impédance du vide. Sa portée provient de la construction parallèle des deux lecteurs sur \(X\).

\subsection{2e confluence : action, gravité et forme électrogravitationnelle}

Soit

\begin{equation}
 Q_P=E_P,
 \qquad
 \ell_P,
 \qquad
 Q_P\ell_P=\hbar c,
 \qquad
 \frac{\ell_P}{Q_P}=\frac{G}{c^4}.
 \label{rm:eq:confluence-planck-pair}
\end{equation}

Le produit est le quantum d’échange \(\hbar c\) ; le quotient est la compliance gravitationnelle. Introduisons

\begin{equation}
 \mathscr T_P=\frac{Q_P}{\ell_P}=\frac{c^4}{G},
 \qquad
 \mathscr C_P=\frac{\ell_P}{Q_P}=\frac{G}{c^4}.
 \label{rm:eq:confluence-tension-compliance}
\end{equation}

L’équation d’Einstein adopte la forme constitutive

\begin{equation}
 G_{\mu\nu}+\Lambda g_{\mu\nu}
 =8\pi\mathscr C_P T_{\mu\nu}.
 \label{rm:eq:confluence-einstein-constitutive}
\end{equation}

Pour une charge \(q\), la coordonnée électrique portée au plan de l’énergie est

\begin{equation}
 \mathcal Q(q)=Q_P\sqrt{\frac{Z}{4\pi\hbar}}\,q.
 \label{rm:eq:confluence-electric-coordinate}
\end{equation}

Newton et Coulomb se réunissent alors dans la forme bilinéaire

\begin{equation}
 \boxed{
 V(r)=\frac{\mathscr C_P}{r}
 \bigl(\mathcal Q_1\mathcal Q_2-E_1E_2\bigr).}
 \label{rm:eq:confluence-newton-coulomb}
\end{equation}

Le cône \(\mathcal Q^2=E^2\) reproduit la condition statique d’extrémalité dans la carte d’Einstein--Maxwell. Le produit des feuillets du vide régit le cône causal ; leur quotient régit le relèvement électrique et l’orientation.

L’identité croisée avec le secteur électrofaible est

\begin{equation}
 \boxed{
 \frac{\hbar^2}{8\pi\ell_P^2}
 \frac{(\kappa/\mu)}{e^2}
 =\frac1{4\pi\alpha_{\HMT}}
 =\frac1{g^2}+\frac1{g'^2}.}
 \label{rm:eq:confluence-einstein-maxwell-ew}
\end{equation}

Cette égalité conserve sa valeur sous l’inversion du feuillet d’action : \(G\) et \(e^2\) se transforment de manière contragrédiente, tandis que \(G\hbar\), \(Ge^2\) et \(\ell_P^2\) restent invariants.

\subsection{3e confluence : géométrie, information et mélange}

Soit \(N=N_{90}+N_{120}\) le nombre total d’excitations de bord et

\begin{equation}
 D_{90/120}=90N_{90}+120N_{120}.
 \label{rm:eq:confluence-modular-provenance}
\end{equation}

L’aire lit \(N\), tandis que le Hamiltonien holographique modulaire lit \(D_{90/120}\). Comme

\begin{equation}
 \det\begin{pmatrix}1&1\\90&120\end{pmatrix}=30,
 \label{rm:eq:confluence-modular-determinant}
\end{equation}

le couple reconstruit exactement

\begin{equation}
 N_{90}=4N-\frac{D_{90/120}}{30},
 \qquad
 N_{120}=\frac{D_{90/120}}{30}-3N.
 \label{rm:eq:confluence-sector-reconstruction}
\end{equation}

La géométrie totale et l’énergie modulaire pondérée conservent donc la quantité et la provenance.

L’état angulaire bicouche

\begin{equation}
 \rho_y=\frac{e^{-yR}}{2\cosh y}
 \label{rm:eq:confluence-rhoy}
\end{equation}

possède une entropie paire et une divergence orientée

\begin{equation}
 S(\rho_y)=\log(2\cosh y)-y\tanh y,
 \qquad
 D(\rho_y\Vert\rho_{-y})=2y\tanh y.
 \label{rm:eq:confluence-info-orientation}
\end{equation}

L’involution \(y\mapsto-y\) est représentée dans le secteur de saveur par

\begin{equation}
 R_{\rm CKM}
 =M_{\rm CKM}\operatorname{diag}(1,-1,1)M_{\rm CKM}^{-1},
 \qquad
 R_{\rm CKM}^2=I,
 \qquad
 \det R_{\rm CKM}=-1.
 \label{rm:eq:confluence-rckm}
\end{equation}

Barbero lit la composante impaire de la transition \(6\to8\), l’entropie lit sa distribution paire et \(R_{\rm CKM}\) réalise la même inversion dans le codomaine de mélange. Ce sont des représentations coordonnées d’une charnière commune, dotées de domaines et d’observables propres.

\subsection{4e confluence : écran, torsion et mémoire hélicoïdale}

Dans le quotient lorentzien \(Q_{\square}\), la soudure et la réponse positive satisfont

\begin{equation}
 \boxed{
 (\beta_m^{\square})^*\Lambda_m^{\square}\beta_m^{\square}
 =\frac{56}{81}I_{Q_{\square}}.}
 \label{rm:eq:confluence-screen-energy}
\end{equation}

Dans le registre profond, l’opérateur de vis vérifie

\begin{equation}
 \boxed{
 V^{-1}\mathscr SV=(2+\sqrt3)\ii\,\mathscr S.}
 \label{rm:eq:confluence-screw}
\end{equation}

La première identité mesure le coût minimal de réponse de l’écran ; la seconde transporte l’orientation et la profondeur. Catalan apparaît comme le deuxième moment de la composante quaternaire. Si \(N_{\rm dep}\delta_n=n\delta_n\) désigne l’opérateur de profondeur du registre hélicoïdal, alors

\begin{equation}
 G_{\rm Cat}=\Tr(N_{\rm dep}^{-2}\mathcal Q_{\chi,-4}).
 \label{rm:eq:confluence-catalan}
\end{equation}

La loi de Perron

\begin{equation}
 M_n=\left(\frac{a_n}{a_0}\right)^{-6}
 \label{rm:eq:confluence-sixth-law}
\end{equation}

fournit la covariance exacte d’une densité quadratique de spin. Le réalisateur de Cartan doit transporter \cref{rm:eq:confluence-screen-energy,rm:eq:confluence-screw} vers un courant et une contorsion dont le mode homogène satisfait \cref{rm:eq:confluence-sixth-law}. Cet enchaînement situe l’interface physique dans une flèche unique et falsifiable.

\subsection{Confluence entre périodicité de l’horloge, énergie d’horizon et coût informationnel}

L’horloge \(108\) fixe

\begin{equation}
 E_P=k_B\Theta_{\rm clk}\log3.
 \label{rm:eq:confluence-planck-trit}
\end{equation}

Pour la région de Sitter de rayon \(R=\ell_P\), l’énergie du vide, la température et l’entropie satisfont

\begin{equation}
 E_\Lambda=T_H S_H=\frac{E_P}{2}
 =\frac{k_B\Theta_{\rm clk}\log3}{2}.
 \label{rm:eq:confluence-horizon-half}
\end{equation}

Il en résulte

\begin{equation}
 E_\Lambda t_P=\frac{\hbar}{2},
 \qquad
 E_\Lambda T_{108}=\frac{h}{2},
 \qquad
 \frac{S_H}{k_B}=\pi.
 \label{rm:eq:confluence-half-action}
\end{equation}

Le volume du vide, l’aire de l’horizon, la température du cycle et le coût informationnel d’un trit expriment la même énergie. La multiplicité \(e^{S_H/k_B}=e^\pi\) coïncide exactement avec la valeur propre expansive de \(e^{\pi J_\varphi}\) ; promouvoir cette égalité en équivalence de microétats exige un entrelaceur préservant la mesure et la dynamique.

\subsection{Confluence fibrée des réponses HMT–MD : vide, action, gravité, information, mélange, mémoire et cosmologie}

\begin{theorem}[Confluence des réponses HMT--MD]
L’état enrichi \(X\) et les primitives déclarées dans chaque lecteur étant fixés, les publications de vide, d’action, de gravité, d’information, de mélange, d’écran, de mémoire hélicoïdale et de cosmologie sont les projections ou compositions définies sur le produit fibré de graphes \cref{rm:eq:master-fiber-product}. Leurs relations croisées sont engendrées par \cref{rm:eq:confluence-ew-master,rm:eq:confluence-newton-coulomb,%
rm:eq:confluence-einstein-maxwell-ew,rm:eq:confluence-sector-reconstruction,%
rm:eq:confluence-info-orientation,rm:eq:confluence-screen-energy,%
rm:eq:confluence-screw,rm:eq:confluence-horizon-half}.
Chaque évaluation numérique est terminale et chaque mémoire orientée demeure dans le facteur qui l’engendre.
\end{theorem}

\begin{proof}
La construction de \(\mathfrak M_X\) fixe l’antécédent commun. Les équations du vide résultent du calcul fonctionnel de \(T\) ; les équations électrofaibles résultent de \(D\) et \(Q(D)\) ; la forme électrogravitationnelle résulte des deux identités du couple de Planck ; la reconstruction \(90/120\) résulte de l’inversion d’une matrice entière de déterminant \(30\) ; l’orientation informationnelle résulte du logarithme de \(\rho_y\) ; l’énergie d’écran résulte de la réponse unique dans \(Q_{\square}\) ; la vis résulte des deux covariances de Weyl et de Pell ; et l’égalité d’horizon résulte des formules géométriques de Sitter spécialisées en \(R=\ell_P\). Aucune preuve n’utilise la décimale terminale pour sélectionner la branche. Le théorème d’indépendance horizontale interdit de transformer les relations croisées en factorisations non démontrées.
\end{proof}

\subsection{Corollaire de confluence métrologique entre les réalisations HMT et SI}

Toute sortie admet le quintuplet indissociable

\begin{equation}
 \mathcal K(P)=
 \bigl(
 \text{antécédent},
 \text{équation},
 \text{invariant},
 \text{signification},
 \text{valeur terminale}
 \bigr).
 \label{rm:eq:confluence-quintuple}
\end{equation}

L’atlas métrologique du volume énumère ces quintuplets. Son contenu probatoire est le suivant : l’antécédent fixe la généalogie ; l’équation fixe l’opération ; l’invariant démontre ce qui subsiste ; la signification explicite la lecture physique ; la valeur terminale permet la reconnaissance externe. Séparer une colonne détruit l’information causale ; conserver les cinq permet d’insérer le résultat dans une carte d’unités sans faire de cette carte un fondement.

\begin{proofstatusbox}
Le produit fibré, les identités algébriques et le théorème d’indépendance horizontale constituent une formalisation exacte. Les équations internes citées conservent le statut démontré dans les chapitres qui contiennent leurs démonstrations complètes. Les lectures d’Einstein--Maxwell, de Sitter, d’Einstein--Cartan et électrofaible demeurent dans leurs interfaces déclarées lorsqu’elles nécessitent un réalisateur supplémentaire. Le théorème organise ces couches sans les promouvoir par similitude numérique.
\end{proofstatusbox}

\begin{falsifierbox}
La confluence est réfutée si deux composantes attribuées au même élément de \(\mathfrak M_X\) proviennent en réalité d’états incompatibles, si l’une des identités génératrices échoue, si une évaluation terminale intervient dans la sélection de son propre antécédent, ou si une mémoire distinguée par un autre lecteur est perdue. Une prétendue factorisation horizontale est réfutée par toute paire d’états ayant le même lecteur intermédiaire et des sorties finales différentes. Les promotions physiques sont en outre confrontées à leurs propres réalisateurs, identités de jauge, de Bianchi, de raffinement et conditions aux limites.
\end{falsifierbox}
